% Draft short note: proofs of the two conjectured formulas in OEIS A334504.
% Status: DRAFT for internal review (not submitted anywhere). See paper/README.md.
\documentclass[11pt]{article}
\usepackage[margin=1in]{geometry}
\usepackage{amsmath,amssymb,amsthm}
\usepackage[hidelinks]{hyperref}

\newtheorem{theorem}{Theorem}
\newtheorem{lemma}[theorem]{Lemma}
\newtheorem{corollary}[theorem]{Corollary}
\theoremstyle{definition}
\newtheorem{definition}[theorem]{Definition}
\theoremstyle{remark}
\newtheorem{remark}[theorem]{Remark}

\newcommand{\Z}{\mathbb{Z}}
\newcommand{\F}{\mathbb{F}}
\newcommand{\lb}{\lfloor \log_2 w \rfloor}

\title{The period of a single cell under elementary cellular automaton rule~26\\ on a cyclic lattice}
\author{[Author names to be decided]}
\date{Draft, November 2026}

\begin{document}
\maketitle

\begin{abstract}
Let $a(w)$ be the eventual period of the orbit of a single live cell under
elementary cellular automaton rule~26 on a cyclic lattice of $w$ cells
(OEIS A334504). Two formulas for $a$ were conjectured in that entry in 2024:
$a(2n+1) = (2n+1)\,2^{A070939(n)}$ and $a(2n+2) = A268754(n)$ for $n \ge 1$.
We prove both. For odd $w \ge 3$ the orbit is purely periodic with period
exactly $w \cdot 2^{\lfloor \log_2 w \rfloor}$; the proof describes every state
explicitly as a window of Pascal's triangle modulo~$2$ drifting around the
ring. For even $w = 2m \ge 4$ we show that from time~$1$ on, the orbit is a
mirror-symmetric embedding of the orbit of a single end cell under rule~90 on
a segment of $m-1$ cells with null boundaries, which is the system whose
period A268754 records. Both results have been checked in the Lean~4 proof
assistant.
\end{abstract}

\section{Introduction}

Rule~26 is the elementary cellular automaton whose new cell value is~$1$
exactly when the neighbourhood (left, centre, right) is $100$, $011$ or $001$.
It is not additive, but started from a single live cell on the infinite line
it draws the same Sierpi\'nski pattern as the additive rule~90, a fact that is
well known (see e.g.\ Wolfram's \emph{A New Kind of Science}). On a finite
cyclic lattice the two rules differ, because the pattern wraps around and
meets itself.

OEIS A334504 lists the eventual period $a(w)$ of a single cell under rule~26
on a cyclic lattice of width~$w$. Its FORMULA field records two conjectures
(A.~Zabolotskiy, 2024): for $n \ge 1$,
\begin{equation}\label{eq:conj}
a(2n+1) = (2n+1)\,2^{A070939(n)}, \qquad a(2n+2) = A268754(n),
\end{equation}
where $A070939(n)$ is the number of binary digits of~$n$ and $A268754(n)$ is
the period of the $n \times 1$ oscillator in the Life-like rule B1/S, whose
OEIS program iterates rule~90 on $n$ cells with null boundaries from a single
end cell. The cross-reference in A268754 likewise only says that it
``seems to be'' a shifted bisection of A334504. We have not found a proof of
either formula in the literature we searched (OEIS, arXiv, OpenAlex; see
Section~\ref{sec:discussion}).

Since $A070939(n) = \lfloor \log_2 n \rfloor + 1 = \lfloor \log_2 (2n+1) \rfloor$
for $n \ge 1$, the first formula is the case $w = 2n+1$ of
Theorem~\ref{thm:odd} below; the second is Theorem~\ref{thm:even}.

\section{Definitions}

Fix $w \ge 1$. A state is a word $x = (x(0), \dots, x(w-1)) \in \F_2^{w}$,
indices taken modulo~$w$; the left neighbour of cell~$i$ is cell $i-1$ and the
right neighbour is cell $i+1$. Write
\[
f(l,c,r) = 1 \iff (l,c,r) \in \{100,\,011,\,001\}
\]
for the local rule of rule~26 ($26 = 00011010_2$, bit $4l+2c+r$), and
$F(x)(i) = f(x(i-1), x(i), x(i+1))$ for the synchronous global map. The seed
is $e_0$, the state with only cell~$0$ alive, and $x_t = F^t(e_0)$.

Since the state space is finite, there are least integers $T \ge 0$ and
$p \ge 1$ with $x_{t+p} = x_t$ for all $t \ge T$; $p$ is the \emph{eventual
period} $a(w)$ and $T$ the \emph{transient}.

\begin{remark}
Reversing the orientation of the ring turns rule~26 into its mirror image,
rule~82, and maps the orbit of $e_0$ bijectively onto the mirrored orbit. The
period is therefore independent of which neighbour is called ``left''.
\end{remark}

We write $b_t(j) = \binom{t}{j} \bmod 2$ (with $b_t(j) = 0$ for $j > t$), so
that $b_t(0) = 1$, $b_0(j) = [j = 0]$ and
$b_{t+1}(j+1) = b_t(j) \oplus b_t(j+1)$.

\section{Odd widths}

\begin{theorem}\label{thm:odd}
Let $w \ge 3$ be odd and $k = \lfloor \log_2 w \rfloor$. Then for all $t \ge 0$,
\[
x_t = e_0 \iff w \cdot 2^{k} \mid t .
\]
Consequently the orbit of $e_0$ is purely periodic (transient~$0$) with
period exactly $a(w) = w \cdot 2^{\lfloor \log_2 w \rfloor}$.
\end{theorem}

The proof rests on an explicit description of every state.

\begin{definition}
For $0 \le s < w$ let $Y_t(s) = 1$ if $s$ is even and $b_t(s/2) = 1$, and
$Y_t(s) = 0$ otherwise. We call $s$ a \emph{window coordinate}.
\end{definition}

Geometrically, row~$t$ of the rule-90 pattern on~$\Z$ has live cells at the
positions $r$ with $r \equiv t \pmod 2$ and $b_t((t+r)/2) = 1$; putting
$s = r + t$, $Y_t$ is the part of that row in the window $-t \le r < w - t$.
The window moves one cell to the left per step.

\begin{lemma}[closed form]\label{lem:closed}
Let $w \ge 3$ be odd. For all $t \ge 0$ and $0 \le i < w$,
\[
x_t(i) = Y_t\big((i+t) \bmod w\big).
\]
\end{lemma}

\begin{proof}
Induction on~$t$. For $t = 0$, $Y_0(s) = [s = 0]$, so the right-hand side is
$e_0$. For the inductive step fix $i$ and let $s = (i+t+1) \bmod w$. The cells
$i-1$, $i$, $i+1$ have window coordinates $s-2$, $s-1$, $s$ (mod~$w$) at
time~$t$, so it suffices to show
\[
f\big(Y_t(s-2),\,Y_t(s-1),\,Y_t(s)\big) = Y_{t+1}(s) \qquad (\text{indices mod } w).
\]
Because $w \ge 3$, the values $0,1,2$ are distinct mod~$w$ and the following
four cases are exhaustive.
\begin{enumerate}
\item $s \ge 2$ even, $s = 2j+2$. The arguments are $(b_t(j), 0, b_t(j+1))$.
  With centre~$0$, $f(l,0,r) = l \oplus r$, so the value is
  $b_t(j) \oplus b_t(j+1) = b_{t+1}(j+1) = Y_{t+1}(s)$.
\item $s \ge 3$ odd. Then $s-2$ and $s$ are odd, both outer arguments are~$0$,
  and $f(0,c,0) = 0 = Y_{t+1}(s)$.
\item $s = 1$. The arguments are $(Y_t(w-1), Y_t(0), Y_t(1)) = (l, 1, 0)$, and
  $f(l,1,0) = 0 = Y_{t+1}(1)$. (Here rule~26 differs from rule~90, which maps
  $110$ to~$1$; this is where the wrapped pattern is cut off.)
\item $s = 0$. The arguments are $(Y_t(w-2), Y_t(w-1), Y_t(0))$. Since $w$ is
  odd, $w-2$ is odd and $Y_t(w-2) = 0$; also $Y_t(0) = b_t(0) = 1$. Hence the
  value is $f(0,c,1) = 1 = Y_{t+1}(0)$. \qedhere
\end{enumerate}
\end{proof}

The oddness of~$w$ enters only in case~4: it makes the two ends $s = w-1$ and
$s = 0$ of the window both even, so the left edge of the pattern always reads
$0\,?\,1$ and is reproduced.

\begin{lemma}[Lucas, $p = 2$]\label{lem:lucas}
\begin{enumerate}
\item[(a)] If $2^{m} \mid t$, then $\binom{t}{j}$ is even for $0 < j < 2^{m}$.
\item[(b)] If $2^{m} \nmid t$, then $\binom{t}{2^{v}}$ is odd for some $v < m$.
\end{enumerate}
\end{lemma}

\begin{proof}
Both follow from Lucas' theorem; we give the short self-contained argument
that is also the one formalised. From $(1+X)^{2a} = (1+X^2)^a$ in $\F_2[X]$,
$b_{2a}(2c) = b_a(c)$ and $b_{2a}(2c+1) = 0$; also $b_t(1) = t \bmod 2$.
(a) Induction on~$m$: write $t = 2a$ with $2^{m-1} \mid a$. Odd $j$ give
$b_t(j) = 0$, and $j = 2j'$ gives $b_t(j) = b_a(j')$ with $0 < j' < 2^{m-1}$.
(b) Take $v = v_2(t) < m$: writing $t = 2^{v} u$ with $u$ odd, repeated use of
$b_{2a}(2c) = b_a(c)$ gives $b_t(2^{v}) = b_u(1) = 1$.
\end{proof}

\begin{proof}[Proof of Theorem~\ref{thm:odd}]
Since $w \ge 3$ is odd, $k \ge 1$ and $2^{k} < w < 2^{k+1}$. By
Lemma~\ref{lem:closed}, $x_t = e_0$ if and only if
$Y_t((i+t) \bmod w) = [i = 0]$ for all $0 \le i < w$.

($\Rightarrow$) Suppose $x_t = e_0$. If $r = t \bmod w \ne 0$, then
$i = w - r$ satisfies $0 < i < w$ and $x_t(i) = Y_t(0) = 1$, a contradiction;
so $w \mid t$ and $x_t(i) = Y_t(i)$. If $2^{k} \nmid t$,
Lemma~\ref{lem:lucas}(b) gives $v < k$ with $b_t(2^{v}) = 1$; then
$i = 2^{v+1} \le 2^{k} < w$ is a nonzero cell with $x_t(i) = Y_t(2^{v+1}) = 1$,
again a contradiction. Hence $2^{k} \mid t$.

($\Leftarrow$) If $w \mid t$ and $2^{k} \mid t$, then $x_t(i) = Y_t(i)$.
For $i = 0$ this is~$1$, for odd~$i$ it is~$0$, and for $i = 2j \ne 0$ we have
$0 < j < 2^{k}$ (because $2j < w < 2^{k+1}$), so $b_t(j) = 0$ by
Lemma~\ref{lem:lucas}(a).

As $\gcd(w, 2^{k}) = 1$, the condition $w \mid t \wedge 2^{k} \mid t$ is
$w 2^{k} \mid t$. Since $x_0 = e_0$ lies on the cycle, the transient is~$0$
and the period is the least positive return time $w 2^{k}$.
\end{proof}

\begin{remark}
The width $w = 1$ is excluded. There the single cell sees itself on both
sides, $f(1,1,1) = 0$, so $x_1 = 0$ is a fixed point: the period $a(1) = 1$
happens to equal $w\,2^{\lfloor\log_2 w\rfloor}$, but the transient is~$1$.
The OEIS form $(2n+1)2^{A070939(n)}$ with $n = 0$ would give~$2$, since
$A070939(0) = 1$; this is why the conjecture is stated for $n \ge 1$.
\end{remark}

\section{Even widths}

Let $n \ge 1$ and consider rule~90 on a segment of $n$ cells $0, \dots, n-1$
with null boundaries: $y'(j) = y(j-1) \oplus y(j+1)$, where $y(-1) = y(n) = 0$.
Let $y_t$ be the orbit of the single end cell $y_0 = e_0$. This is exactly the
recurrence computed by the program in OEIS A268754 (state as an integer,
update $((h \ll 1) \oplus (h \gg 1))$ masked to $n$ bits, start $h = 1$), and
A268754$(n)$ is the eventual period of $(y_t)$.

\begin{theorem}\label{thm:even}
Let $w = 2m$ with $m \ge 2$ and $n = m - 1$. Define the embedding
$E : \F_2^{n} \to \F_2^{w}$ by $E(y)(0) = E(y)(m) = 0$,
$E(y)(i) = y(i-1)$ for $0 < i < m$ and $E(y)(i) = y(2m-1-i)$ for $m < i < 2m$.
Then $x_{t+1} = E(y_t)$ for all $t \ge 0$. Consequently, for every $p \ge 1$,
$(x_t)$ is eventually $p$-periodic if and only if $(y_t)$ is, so
$a(2n+2) = A268754(n)$ for all $n \ge 1$, and the transient of $(x_t)$ is one
more than that of $(y_t)$.
\end{theorem}

\begin{proof}
\emph{Parity.} Every live cell of $y_t$ lies at a position $j \equiv t
\pmod 2$: true for $t = 0$, and $y_{t+1}(j)$ depends only on positions
$j \pm 1$.

\emph{One step.} Let $y$ have all live cells of one parity and put $x = E(y)$.
Then $x$ has no two adjacent live cells: on each half adjacent cells
correspond to adjacent positions of~$y$, and cells $0$ and $m$ are dead (the
ring has even length, so wrapping around causes no problem). Hence the
neighbourhood $110$ never occurs in~$x$, and on such states rule~26 agrees
with rule~90: $f(l,c,r) = l \oplus r$ whenever $c = 1 \Rightarrow l = 0$. We
check $F(x) = E(y')$ cell by cell. Cell~$0$ sees $y(0)$ on both sides (cells
$2m-1$ and~$1$), giving~$0$; cell~$m$ likewise sees $y(m-2)$ twice. For
$0 < i < m$ the left neighbour is $y(i-2)$ (or the dead cell~$0$ when $i = 1$)
and the right neighbour is $y(i)$ (or the dead cell~$m$ when $i = m-1$, which
matches $y(n) = 0$), so $F(x)(i) = y'(i-1)$. The right half is the mirror image.

\emph{First step.} $x_1 = F(e_0)$ has live cells exactly $1$ and $2m-1$
(neighbourhoods $100$ and $001$), which is $E(e_0)$ because $m \ge 2$.
Induction gives $x_{t+1} = E(y_t)$.

\emph{Periods.} $E$ is injective (its left half reads back~$y$) and
$E(y)(0) = 0 \ne e_0(0)$, so $x_{s+1} = x_{t+1} \iff y_s = y_t$ and $x_0$ never
recurs. Thus $x_{t+p} = x_t$ for all $t \ge T+1$ iff $y_{t+p} = y_t$ for all
$t \ge T$, which gives the claims.
\end{proof}

\section{Discussion}\label{sec:discussion}

\paragraph{What is new and what is not.}
The formulas~\eqref{eq:conj} are not ours; they were conjectured in OEIS
A334504. The ingredients of the proofs are classical: the Sierpi\'nski pattern
of rule~26 on the line, Lucas' theorem, the periods of rule~90, and the
reflection trick relating null and cyclic boundaries. The device of reducing a
non-additive rule to rule~90 on states with parity-separated support also has
precedents (Martin, Odlyzko and Wolfram treat rule~18 this way, and a 2019
comment in OEIS A160657 reduces oscillators of a Life-like rule to finite
rule-90 wires).% TODO(verify): read MOW 1984 Sec. on rule 18 and the A160657 comment before release.
What we contribute is the specific analysis of rule~26 on the ring: the
drifting-window closed form (Lemma~\ref{lem:closed}) with its boundary
analysis, the exact return-time characterisation for odd widths, and the
explicit orbit embedding for even widths. We do not claim more.

\paragraph{Literature search.}
We searched OEIS (A334504, A268754 and sequence searches), arXiv and OpenAlex
in November 2026 and found no proof of either formula. MathSciNet, zbMATH,
Google Scholar and the math-fun mailing list where the sequence originated
have not yet been searched; a proof may exist that was not reported to OEIS.

\paragraph{Formal verification.}
Both theorems are formalised in Lean~4 (core library only, no
\texttt{sorry}, no \texttt{native\_decide}; axioms limited to
\texttt{propext}, \texttt{Classical.choice}, \texttt{Quot.sound}) in the
files \texttt{lean/Discovery/Rule26.lean} and
\texttt{lean/Discovery/Rule26Even.lean} of the accompanying repository.
States are encoded as natural numbers whose bit~$i$ is cell~$i$. Two points
of the correspondence are not themselves formal: (i) Lean defines $b_t(j)$ by
Pascal's rule and does not state its equality with $\binom{t}{j} \bmod 2$
(the main theorems do not mention $b$); (ii) the identification of the Lean
definitions with the OEIS definitions is by inspection and numerical
cross-checks (independent simulations agree with the formulas for odd
$w \le 129$ and with A268754 for $n \le 40$). The audit is recorded in
\texttt{paper/spec\_audit.md}.

\paragraph{Use of AI.}
The conjectures were selected, the proofs found, the Lean formalisation
written and the numerical checks run with substantial help from AI systems
(Claude Code, with independent AI review passes); the human authors are
responsible for the final text. A detailed log is in the repository.
% TODO: authors to state their own contribution precisely before any release.

\begin{thebibliography}{9}
\bibitem{oeis334504} OEIS Foundation, Sequence A334504, \emph{The On-Line
  Encyclopedia of Integer Sequences}, \url{https://oeis.org/A334504}.
\bibitem{oeis268754} OEIS Foundation, Sequence A268754,
  \url{https://oeis.org/A268754}.
\bibitem{oeis160657} OEIS Foundation, Sequence A160657,
  \url{https://oeis.org/A160657}.
\bibitem{mow} O.~Martin, A.~M.~Odlyzko, S.~Wolfram, Algebraic properties of
  cellular automata, \emph{Comm. Math. Phys.} 93 (1984), 219--258.
\bibitem{nks} S.~Wolfram, \emph{A New Kind of Science}, Wolfram Media, 2002.
\end{thebibliography}

\end{document}
